% ---
% school: "college"
% level: "8eme"
% section: null
% chapter: "العمليات والحساب على الأعداد الكسرية"
% source: "Axiomaths"
% title: "تفكيك $\dfrac{1"
% country: "Tunisie"
% year: 2026
% professor: "Axiomaths"
% difficulty: "Difficile"
% points: null
% tags: ["fractions", "fractions-egyptiennes", "recherche-exhaustive"]
% ---
% Exercice 15/18 de la série 8ème « العمليات والحساب على الأعداد الكسرية » (statement + \axsol + solution). Compiler avec XeLaTeX + axiomath.sty (+ axiomaths-logo.png).
\documentclass[12pt,a4paper]{article}
\usepackage[level=8b,color=15803D]{axiomath}
\axsetup{title={العمليات والحساب على الأعداد الكسرية}, corrige}

\begin{document}

\begin{corr}[title={تفكيك $\dfrac{1}{6}$}, stars=4, source={Axiomaths}]
نبحث عن عددين صحيحين طبيعيين غير معدومين $a$ و $b$ حيث $a<b$ و $\dfrac{1}{6}=\dfrac{1}{a}+\dfrac{1}{b}$.
\begin{enumerate}
  \item تحقق أن $a=7$ و $b=42$ حل.
  \item بيّن أن $a>6$.
  \item بيّن أن $\dfrac{1}{12}<\dfrac{1}{a}$ ثم استنتج أن $a<12$. [قارن $\dfrac{1}{b}$ و $\dfrac{1}{a}$ ثم قارن $\dfrac{1}{6}$ و $\dfrac{2}{a}$]
  \item من أجل كل قيمة ممكنة لـ $a$ احسب $\dfrac{1}{6}-\dfrac{1}{a}$ وحدّد هل يمكن كتابته على شكل $\dfrac{1}{b}$ حيث $b$ عدد صحيح طبيعي.
  \item استنتج جميع الحلول.
\end{enumerate}
\axsol
\begin{enumerate}[itemsep=14pt]
  \item
  \[\dfrac{1}{7}+\dfrac{1}{42}=\dfrac{6}{42}+\dfrac{1}{42}=\dfrac{7}{42}=\dfrac{1}{6}\]
  \item $\dfrac{1}{6}=\dfrac{1}{a}+\dfrac{1}{b}$ و $\dfrac{1}{b}>0$ ، إذن $\dfrac{1}{a}<\dfrac{1}{6}$.
  من بين عددين موجبين، الأصغر مقلوبا هو الأكبر : إذن $a>6$.
  \item بما أن $a<b$ فإن $\dfrac{1}{b}<\dfrac{1}{a}$. إذن :
  \[\dfrac{1}{6}=\dfrac{1}{a}+\dfrac{1}{b}<\dfrac{1}{a}+\dfrac{1}{a}=\dfrac{2}{a}\]
  نقسم على $2$ :
  \[\dfrac{1}{12}<\dfrac{1}{a}\]
  وبما أن $\dfrac{1}{12}<\dfrac{1}{a}$ فإن $a<12$.
  \item من السؤالين 2 و3 : $a\in\{7;\,8;\,9;\,10;\,11\}$.
  \[a=7:\quad\dfrac{1}{6}-\dfrac{1}{7}=\dfrac{7}{42}-\dfrac{6}{42}=\dfrac{1}{42}\quad\Rightarrow\ b=42\]
  \[a=8:\quad\dfrac{1}{6}-\dfrac{1}{8}=\dfrac{4}{24}-\dfrac{3}{24}=\dfrac{1}{24}\quad\Rightarrow\ b=24\]
  \[a=9:\quad\dfrac{1}{6}-\dfrac{1}{9}=\dfrac{3}{18}-\dfrac{2}{18}=\dfrac{1}{18}\quad\Rightarrow\ b=18\]
  \[a=10:\quad\dfrac{1}{6}-\dfrac{1}{10}=\dfrac{5}{30}-\dfrac{3}{30}=\dfrac{2}{30}=\dfrac{1}{15}\quad\Rightarrow\ b=15\]
  \[a=11:\quad\dfrac{1}{6}-\dfrac{1}{11}=\dfrac{11}{66}-\dfrac{6}{66}=\dfrac{5}{66}\]
  والعدد $\dfrac{5}{66}$ ليس من الشكل $\dfrac{1}{b}$ (بسطه $5$ ولا يقبل الاختزال) : لا توجد قيمة لـ $b$ في هذه الحالة.
  \item الحلول هي :
  \[\ib{$(a\,;b)\in\{(7\,;42),\ (8\,;24),\ (9\,;18),\ (10\,;15)\}$}\]
  \[\dfrac{1}{6}=\dfrac{1}{7}+\dfrac{1}{42}=\dfrac{1}{8}+\dfrac{1}{24}=\dfrac{1}{9}+\dfrac{1}{18}=\dfrac{1}{10}+\dfrac{1}{15}\]
\end{enumerate}
\textbf{الفكرة :} حصر $a$ بين $6$ و $12$ يحوّل مسألة غير منتهية إلى خمس حالات فقط نختبرها.

\textbf{انتبه :} من أجل $a=12$ نجد $\dfrac{1}{6}-\dfrac{1}{12}=\dfrac{1}{12}$ أي $b=12=a$ ، وهذا مرفوض لأن $a<b$.
\end{corr}

\end{document}
